Es sei
\mathl{\left(   x _n  \right)_{n \in \N }}{} eine
\definitionsverweis {Cauchy-Folge}{}{}
in $X$. Nehmen wir an, dass diese Folge nicht konvergiert. Nach
[[Metrischer Raum/Cauchy-Folge/Häufungspunkt/Konvergent/Aufgabe|Aufgabe]]
besitzt sie dann auch keinen Häufungspunkt. Das bedeutet, dass es zu jedem Punkt

\mavergleichskette
{\vergleichskette
{ y
}
{ \in }{ X
}
{  }{
}
{    }{
}
{   }{
}
}
{}{}{}
eine offene Umgebung

\mavergleichskette
{\vergleichskette
{ y
}
{ \in }{ U_y
}
{  }{
}
{    }{
}
{   }{
}
}
{}{}{}
derart gibt, dass es darin nur endlich viele Folgenglieder gibt. Aufgrund der Kompaktheit gibt es zur Überdeckung

\mavergleichskettedisp
{\vergleichskette
{ X
}
{ =} { \bigcup_{y \in X} U_y
}
{ } {
}
{ } {
}
{ } {
}
}
{}{}{}
eine endliche Teilüberdeckung, also

\mavergleichskettedisp
{\vergleichskette
{ X
}
{ =} { \bigcup_{ i = 1}^m  U_{y_i }
}
{ } {
}
{ } {
}
{ } {
}
}
{}{}{.}
Dann wären ab einem $N$ alle Folgenglieder außerhalb dieser Menge, was absurd ist.

 [[Kategorie:Latexseite]]
